\chapter{Hybrid Congestion Control}
\label{chap:congestion_balancing}
With the user-side job recommender system 
(Chapter~\ref{chap:user_recommendation}) and the recruiter-side 
suitability model (Chapter~\ref{chap:recruiter_recommendation}) 
established, this chapter addresses the congestion introduced in 
Chapter~\ref{chap:intro} that emerges as a consequence of 
recommendation at scale. Throughout this chapter and the congestion 
experiments (Section~\ref{sec:exp_congestion}), each user is assumed 
to apply to exactly their top-$k_u$ recommended jobs, where $k_u$ 
matches their historical application count.

This chapter starts with the job-side market-maker interventions in
Section~\ref{sec:job_congestion} and the candidate-side interventions in
Section~\ref{sec:user_congestion}. It then introduces the bilateral
scoring approach in Section~\ref{sec:bilateral_scoring}, which combines
the user and recruiter models into a single score. Finally, the complete
two-phase pipeline is described in Section~\ref{sec:hybrid_strategy},
integrating all three components into a unified system. Evidence of both
forms of congestion, and the effectiveness of the proposed mitigation
strategies, is evaluated in a simulated environment using the dataset
described in Section~\ref{sec:exp_congestion}.
% ============================================================
\section{Job-Side Market-Maker Intervention}
\label{sec:job_congestion}
% ============================================================

Job congestion arises when the application distribution concentrates applications on a small subset of popular job postings, leaving the majority under-applied. To redistribute traffic more equitably, a market-maker modification is applied to the job ranking score during the recommendation phase. The modification is a function of each job's current application count at the time a user is being processed, so that already-popular jobs are progressively penalised and under-applied jobs receive a visibility boost.

Concretely, at the moment user $u$ is being processed, let $c_j^{(t)}$ denote the number of applications job $j$ has already received from previously processed users in the current simulation run. The modified ranking score for the $(u,j)$ pair is:
\begin{equation}
    \hat{s}_{uj} = s_{uj} + \delta_{\text{job}}\!\left(c_j^{(t)}\right),
\end{equation}
where $s_{uj}$ is the base relevance score from the \gls{TT} model,
$\delta_{\text{job}}(\cdot)$ is a job congestion modification function. Jobs within the safe zone $[\tau_{\min}^{\text{job}}, \tau_{\max}^{\text{job}}]$ receive no
modification ($\delta_{\text{job}} = 0$). Starved jobs ($c_j^{(t)} < \tau_{\min}^{\text{job}}$) receive a positive boost to increase their visibility. Congested jobs ($c_j^{(t)} > \tau_{\max}^{\text{job}}$) receive a negative penalty to reduce
further applications. Jobs where $c_j^{(t)} \geq H^{\text{job}}$ for any time step $t$ are penalised heavily to ensure they are at the bottom of the ranking. Three functional forms for $\delta_{\text{job}}$ are considered, each making a different assumption about how aggressively congestion should be
penalised.

\paragraph{Linear modification.}
The penalty grows at a constant rate per unit of congestion beyond the safe
zone:
\begin{equation}
    \delta_{\text{lin}}(c) = \begin{cases}
        \min\!\bigl(\gamma_b \cdot (\tau_{\min}^{\text{job}} - c),\; B\bigr) & c < \tau_{\min}^{\text{job}}, \\
        0 & \tau_{\min}^{\text{job}} \leq c \leq \tau_{\max}^{\text{job}}, \\
        \max\!\bigl(-\gamma_p \cdot (c - \tau_{\max}^{\text{job}}),\; -P\bigr) & c > \tau_{\max}^{\text{job}},
    \end{cases}
\end{equation}
where $B$ and $P$ are the maximum boost and penalty magnitudes respectively,
and $\gamma_b$, $\gamma_p$ are the rates per unit of deviation from the safe
zone.

\paragraph{Quadratic modification.}
The penalty grows quadratically, starting slowly for mildly congested jobs and
accelerating as congestion increases:
\begin{equation}
    \delta_{\text{quad}}(c) = \begin{cases}
        \min\!\bigl(\kappa_b \cdot (\tau_{\min}^{\text{job}} - c)^2,\; B\bigr) & c < \tau_{\min}^{\text{job}}, \\
        0 & \tau_{\min}^{\text{job}} \leq c \leq \tau_{\max}^{\text{job}}, \\
        \max\!\bigl(-\kappa_p \cdot (c - \tau_{\max}^{\text{job}})^2,\; -P\bigr) & c > \tau_{\max}^{\text{job}},
    \end{cases}
\end{equation}
where $\kappa_b$ and $\kappa_p$ are scaling parameters.

\paragraph{Sigmoid modification.}
A smooth S-curve naturally bounds the modification within $[-P, B]$ without
requiring explicit clamping:
\begin{equation}
    \delta_{\text{sig}}(c) = \begin{cases}
        B \cdot \!\left(\dfrac{2}{1 + e^{-\kappa_s(\tau_{\min}^{\text{job}} - c)}} - 1\right)
            & c < \tau_{\min}^{\text{job}}, \\[6pt]
        0 & \tau_{\min}^{\text{job}} \leq c \leq \tau_{\max}^{\text{job}}, \\[6pt]
        -P \cdot \!\left(\dfrac{2}{1 + e^{-\kappa_s(c - \tau_{\max}^{\text{job}})}} - 1\right)
            & c > \tau_{\max}^{\text{job}},
    \end{cases}
\end{equation}
where $\kappa_s$ controls the steepness of the curve.

All three functions share the same boundary values ($B$, $P$,
$H^{\text{job}}$) but differ in how quickly the modification grows
between the safe zone and the hard cap. The linear form applies a
constant rate of change, reaching its maximum penalty only at the hard
cap. The quadratic form is lenient on moderately congested jobs but
accelerates once congestion becomes severe, concentrating the penalty
on the heaviest tail. The sigmoid form is the most aggressive early
on, applying a steep penalty as soon as a job crosses the congested
threshold $\tau_{\max}^{\text{job}}$, before levelling off toward the hard cap. The specific parameter values of the boundary values and the behaviour of the modification functions used in the experiments are reported in Section~\ref{sec:exp_congestion_setting}. The empirical results of the job-side market-maker intervention are presented in Section~\ref{sec:exp_job_congestion}.
% ============================================================
\section{Candidate-Side Market-Maker Intervention}
\label{sec:user_congestion}
% ============================================================

The candidate congestion arises when recruiter attention concentrates on a small
subset of candidates: strong applicants accumulate shortlisting events across
many job inboxes simultaneously, while weaker-signaled but potentially
qualified candidates are overlooked entirely. This congestion is strongly correlated with job congestion as popular jobs attract more applications, creating larger inboxes and more competition for recruiter attention. To mitigate this, the same market-maker framework
developed for the job side is adapted here to act on the recruiter's shortlisting decision rather than the user's application decision.

At the moment job $j$'s inbox is being processed, let $v_u^{(t)}$ denote the
number of times candidate $u$ has already been shortlisted in the current
simulation run. The modified recruiter ranking score is:
\begin{equation}
    \hat{r}_{uj} = r_{uj} + \delta_{\text{cand}}\!\left(v_u^{(t)}\right),
\end{equation}
where $r_{uj}$ is the recruiter model's base suitability probability and
$\delta_{\text{cand}}(\cdot)$ --- the candidate-side congestion modification function --- uses the same three functional forms as
Section~\ref{sec:job_congestion} but parameterised with candidate-side
thresholds $\tau_{\min}^{\text{cand}}$ and $\tau_{\max}^{\text{cand}}$. The same values for ($B$, $P$) are used and a candidate-side hard cap $H^{\text{cand}}$ is applied. The top-$k_{\text{short}}$ candidates by $\hat{r}_{uj}$ are shortlisted for each job,
where $k_{\text{short}}$ is the fixed shortlist size per job. As more
job inboxes are processed sequentially, candidates who have been shortlisted
frequently accumulate negative modifications, naturally redirecting recruiter
attention toward candidates who have not yet been surfaced. The candidate-side
parameter values and the behaviour of the modification functions used in the experiments are reported in
Section~\ref{sec:exp_congestion_setting} and the empirical results are presented in Section~\ref{sec:exp_user_congestion}.

% ============================================================
\section{Bilateral Scoring}
\label{sec:bilateral_scoring}
% ============================================================

The job-side market-maker from Section~\ref{sec:job_congestion} reduces 
congestion by modifying ranking scores based on application counts, but 
it does not address a more fundamental issue: the \gls{TT} model ranks 
jobs purely from the user's perspective, blind to whether the user is 
actually a competitive candidate for the role. A user may be directed 
toward a popular job because it matches their profile, yet have little 
chance of being shortlisted by the recruiter. This creates a mismatch 
between what is recommended and what is actionable.

To address this, the recruiter-side suitability model is incorporated 
directly into the job ranking score. The \gls{TT} logit 
$\hat{y}_{\text{TT}}$ is first transformed to a probability 
$p_{\text{TT}} = \sigma(\hat{y}_{\text{TT}}) \in [0,1]$, and the 
XGBoost recruiter model produces a shortlisting probability 
$p_{\text{rec}} \in [0,1]$ for the same pair. The bilateral ranking 
score for user $u$ on job $j$ is then:
\begin{equation}
    s_{uj} = \alpha \cdot p_{\text{TT}} + (1 - \alpha) \cdot p_{\text{rec}},
    \label{eq:bilateral_score}
\end{equation}
where $\alpha \in [0,1]$ controls the balance between user preference 
and recruiter suitability. Both signals are on the same $[0,1]$ 
probability scale, making the weighted average well-defined and 
interpretable. The effect of this 
implicit redistribution is examined in 
Section~\ref{sec:exp_bilateral}.

% ============================================================
\section{Complete Two-Phase Pipeline}
\label{sec:hybrid_strategy}
% ============================================================

Building on the bilateral scoring from
Section~\ref{sec:bilateral_scoring} and the market-maker interventions
from Sections~\ref{sec:job_congestion} and~\ref{sec:user_congestion},
this section combines all three components into a unified two-phase
pipeline. The bilateral score ensures that recommended jobs are
actionable from both sides of the market; the market-maker corrections
then redistribute attention where the scoring alone is insufficient to
prevent congestion.

Because the system is two-sided, there is no single recommendation
list --- each phase produces a different output for a different audience:

\begin{itemize}
    \item \textbf{Phase~1} answers the question \emph{``which jobs
    should user $u$ apply to?''} and produces a ranked job list per
    user. A simulated application pool is generated by taking the
    top-$k_u$ jobs for each user, where $k_u$ matches their historical
    application count.
    \item \textbf{Phase~2} answers the question \emph{``among all
    applicants to job $j$, who should be shortlisted?''} and produces
    a candidate ranking per job, taking the application pool from
    Phase~1 as its input. The top-$k_{\text{short}}$ candidates by the modified recruiter score are shortlisted for each job.
\end{itemize}

The two phases use independent scoring formulas and independent running
counters, as they operate in opposite directions on different questions.
Phase~1 decides who is in each job's inbox; Phase~2 decides who gets
surfaced from it. 

\subsection{Phase 1: Job Recommendation with Bilateral Scoring and Job Congestion Correction}
\label{sec:hybrid_job_congestion}

Phase~1 extends the bilateral score from
Equation~\ref{eq:bilateral_score} with the job-side congestion
modification from Section~\ref{sec:job_congestion}:
\begin{equation}
    s_{uj} = \alpha \cdot \bigl(p_{\text{TT}} +
    \delta_{\text{job}}\!\left(c_j^{(t)}\right)\bigr)
    + (1 - \alpha) \cdot p_{\text{rec}},
    \label{eq:combined_score}
\end{equation}
where $\delta_{\text{job}}(\cdot)$ is the job congestion modification,
$c_j^{(t)}$ is the running application count for job $j$ at the time
user $u$ is processed, $p_{\text{TT}}$ and $p_{\text{rec}}$ are the \gls{TT} and recruiter probabilities, respectively, and $\alpha = 0.5$ is used in the experiments.
The job congestion penalty is applied inside the \gls{TT} component
only, as applications flow in the user-to-job direction. The recruiter
probability is left unmodified, since candidate fit should be
independent of how many other users happen to be applying to the same
job. 

\subsection{Phase 2: Recruiter Shortlisting with Candidate Congestion Correction}
\label{sec:hybrid_user_congestion}

Phase~2 takes the job inboxes produced by Phase~1 and determines for each job, which
candidates should the recruiter shortlist. The recruiter suitability
score $p_{\text{rec}}$ is reused from Phase~1, but now ranks
\emph{candidates within a job's inbox} rather than jobs for a user.
The candidate-side market-maker modification from
Section~\ref{sec:user_congestion} is applied at this stage:
\begin{equation}
    \hat{r}_{uj} = p_{\text{rec}} +
    \delta_{\text{cand}}\!\left(v_u^{(t)}\right),
\end{equation}
where $v_u^{(t)}$ is the running count of how many times candidate $u$
has already been shortlisted. The experimental results of the complete two-phase pipeline are presented in Section~\ref{sec:exp_hybrid_results}.